\section{入力短縮と作成費用の関係}
全487保存要求の２表示，974入力を計数し，欠測と不一致は０だった．使用量既知486要求では元の実測入力に一致した．不明１要求の入力も数えられたが，元の終端使用量は不明のまま保持する．計数本体は37.420秒，モデル生成，API，評価の追加は０だった．

\begin{table}[t]
\centering\small
\caption{要約採用後の保存入力差．$D-C$は実測総節約ではない．}
\label{tab:inputs}
\begin{tabular}{llrrr}
\hline
方式 & 課題 & $D$ & 別要約$C$ & $D-C$ \\
\hline
別要約 & 写真 & 15,434 & 18,455 & $-3,021$ \\
別要約 & CSV & 200,507 & 21,419 & 179,088 \\
必須欄 & 写真 & 8,023 & -- & -- \\
必須欄 & CSV & 163,335 & -- & -- \\
\hline
\end{tabular}
\par\smallskip\footnotesize 別要約の値は３群でそれぞれ同じ．必須欄は群3．
同時出力の要約だけの出力費用は未分離であり，--を０とはしない．
\end{table}

表\ref{tab:inputs}の写真の別要約では，要約後に２要求があり，各入力差7,717の合計は作成費用に3,021足りなかった．CSVの別要約では３要求目で累積入力差が31,659となり，作成費用21,419を上回った．終了までの19要求の差は200,507だった．全群の精算と返金及び群2の任意欄は要約未採用で，入力差０だった．

必須欄の写真では要約後の１要求で8,023，CSVでは15要求で163,335の入力差があった．後者で入力短縮が実際に起きても，元の未達は変わらない．CSV別要約の実測440,461に$D$を足した640,968，必須欄438,073に$D$を足した601,408は仕事予算450,000を超える．これらを実行可能な全保持の比較結果としない．復元入力と実際の出力予約がcontext上限を超える例はなかった．
